
\documentclass{article}
\usepackage{colortbl}
\usepackage{makecell}
\usepackage{multirow}
\usepackage{supertabular}

\begin{document}

\newcounter{utterance}

\centering \large Interaction Transcript for game `ifeval', experiment `ifeval', episode 73 with Qwen3.6{-}27B.
\vspace{24pt}

{ \footnotesize  \setcounter{utterance}{1}
\setlength{\tabcolsep}{0pt}
\begin{supertabular}{c@{$\;$}|p{.15\linewidth}@{}p{.15\linewidth}p{.15\linewidth}p{.15\linewidth}p{.15\linewidth}p{.15\linewidth}}
    \# & \multicolumn{2}{c}{Player} && \multicolumn{2}{c}{Game Master} \\
    \hline

    \theutterance \stepcounter{utterance}  
    & & & \multicolumn{4}{p{0.6\linewidth}}{
        \cellcolor[rgb]{0.9,0.9,0.9}{
            \makecell[{{p{\linewidth}}}]{
                \texttt{\tiny{[P1$\langle$GM]}}
                \texttt{Write a planning doc for a software engineer task. Follow the exact format below:} \\
\\ 
\texttt{Part 1. Task Summary} \\
\texttt{{[}put details here{]}} \\
\texttt{***} \\
\texttt{Part 2. Motivation} \\
\texttt{{[}put details here{]}} \\
\texttt{***} \\
\texttt{Part 3. Milestone objectives} \\
\texttt{{[}put details here{]}} \\
\texttt{***} \\
\texttt{Part 4. Timeline} \\
\texttt{{[}put details here{]}} \\
\texttt{***} \\
\texttt{Part 5. Doc history} \\
\texttt{{[}put details here{]}} \\
            }
        }
    }
     \\ \\

    \theutterance \stepcounter{utterance}  
    & \multicolumn{4}{p{0.6\linewidth}}{
        \cellcolor[rgb]{0.561,0.737,0.561}{
            \makecell[{{p{\linewidth}}}]{
                \texttt{\tiny{[P1$\rangle$GM]}}
                \texttt{Part 1. Task Summary} \\
\texttt{**Task Title:** Implementation of Asynchronous Email Notification Service} \\
\texttt{**Owner:** Backend Engineering Team} \\
\texttt{**Priority:** High} \\
\\ 
\texttt{This task involves designing and implementing a new asynchronous email notification service to replace the current synchronous email sending logic within the user registration and password reset flows. The new service will utilize a message queue (RabbitMQ) to decouple the core application logic from email delivery, ensuring that transient failures in the email provider do not block user{-}facing operations. The implementation includes creating a producer service, a consumer worker, and a retry mechanism with exponential backoff.} \\
\\ 
\texttt{***} \\
\texttt{Part 2. Motivation} \\
\texttt{Currently, the application sends emails synchronously via the SMTP provider. This approach has led to several critical issues:} \\
\texttt{1.  **Latency:** User registration and password reset requests take 2{-}4 seconds to complete due to network latency with the email provider, resulting in a poor user experience.} \\
\texttt{2.  **Reliability:** If the email provider experiences an outage, the entire registration flow fails, leading to lost conversions and increased support tickets.} \\
\texttt{3.  **Scalability:** During peak traffic, the synchronous calls consume significant thread pool resources, causing bottlenecks in other API endpoints.} \\
\\ 
\texttt{By moving to an asynchronous model, we aim to reduce API response times to under 200ms, improve system resilience against email provider outages, and free up application resources for core business logic.} \\
\\ 
\texttt{***} \\
\texttt{Part 3. Milestone objectives} \\
\texttt{1.  **Design \& Schema Definition:**} \\
\texttt{    *   Define the message schema for email notifications (JSON format).} \\
\texttt{    *   Design the database schema for the `notification\_logs` table to track delivery status.} \\
\texttt{    *   Finalize the RabbitMQ exchange and queue topology.} \\
\\ 
\texttt{2.  **Core Implementation:**} \\
\texttt{    *   Implement the `EmailProducer` service to publish messages to the queue.} \\
\texttt{    *   Implement the `EmailConsumer` worker to fetch messages and send emails via the SMTP provider.} \\
\texttt{    *   Integrate the retry logic with exponential backoff (max 3 retries).} \\
\\ 
\texttt{3.  **Integration \& Testing:**} \\
\texttt{    *   Update the User Registration and Password Reset controllers to use the new producer.} \\
\texttt{    *   Write unit tests for producer/consumer logic.} \\
\texttt{    *   Write integration tests to verify end{-}to{-}end flow using a local RabbitMQ instance.} \\
\texttt{    *   Perform load testing to verify latency improvements.} \\
\\ 
\texttt{4.  **Deployment \& Monitoring:**} \\
\texttt{    *   Deploy the consumer worker as a separate microservice/container.} \\
\texttt{    *   Set up alerts for queue depth and failed delivery rates.} \\
\texttt{    *   Gradually roll out the change using feature flags.} \\
\\ 
\texttt{***} \\
\texttt{Part 4. Timeline} \\
\texttt{*   **Week 1:** Design phase, schema definition, and infrastructure setup (RabbitMQ configuration).} \\
\texttt{*   **Week 2:** Development of Producer and Consumer services; implementation of retry logic.} \\
\texttt{*   **Week 3:** Integration with existing controllers; comprehensive unit and integration testing.} \\
\texttt{*   **Week 4:** Code review, staging environment deployment, load testing, and final production rollout.} \\
\\ 
\texttt{***} \\
\texttt{Part 5. Doc history} \\
\texttt{| Date | Author | Changes |} \\
\texttt{| :{-}{-}{-} | :{-}{-}{-} | :{-}{-}{-} |} \\
\texttt{| 2023{-}10{-}25 | Jane Doe | Initial draft created. |} \\
\texttt{| 2023{-}10{-}26 | John Smith | Added specific retry logic requirements and monitoring metrics. |} \\
\texttt{| 2023{-}10{-}27 | Jane Doe | Updated timeline to reflect additional testing phase. |} \\
            }
        }
    }
    & & \\ \\

    \theutterance \stepcounter{utterance}  
    & & & \multicolumn{2}{p{0.3\linewidth}}{
        \cellcolor[rgb]{0.9,0.9,0.9}{
            \makecell[{{p{\linewidth}}}]{
                \texttt{\tiny{[GM$|$GM]}}
                {'length_constraints:number_paragraphs': {'num_paragraphs': 5}}
            }
        }
    }
    & & \\ \\

    \theutterance \stepcounter{utterance}  
    & & & \multicolumn{2}{p{0.3\linewidth}}{
        \cellcolor[rgb]{0.9,0.9,0.9}{
            \makecell[{{p{\linewidth}}}]{
                \texttt{\tiny{[GM$|$GM]}}
                \texttt{game\_result = WIN} \\
            }
        }
    }
    & & \\ \\

\end{supertabular}
}

\end{document}
